\section{Limitations and threats to validity}\label{sec:limits}

\paragraph{The comparisons are descriptive, not controlled.}
The stages of \cref{tab:stages} differ at once in seed origin, walk length, format range and date, so their yields cannot be attributed to one factor. We did not measure whether directed sampling raises the probability of finding a lower rank compared with uniform sampling, nor whether a longer walk beats several short ones of equal total cost, nor whether the basis rewrite is necessary for any particular record (a later walk might reach the same rank from the plain projection). Every such statement in this paper is a statement about what happened.

\paragraph{Randomness.}
Walk streams are offsets on one full-period linear congruential cycle selected by small, correlated integer nonces, and all seeds of one format within one log share a nonce. We did not test the independence of the streams. Determinism across platforms and compiler versions is untested; only same-machine rebuilds of the same source were replayed, and the binary that ran the first three days was lost with the scratch directory (the sources were unchanged and all 11 first-day walks replay byte for byte with the later build).

\paragraph{Verification.}
Both exactness implementations were written by the same author and share the bit-layout convention. Only 42 of the 244 distinct walks were replayed (\cref{sec:lineage}); closure rows are not rebuilt from leaves by the verifier; five replayed seeds have incomplete leaf provenance; three intermediate witnesses could not be recovered. Exactness of every claimed tensor is established directly and is not affected by any of these.

\paragraph{The campaign had defects that we found and that remain in the data.}
The self-referential child selection (\cref{sec:selfref}) left at least 48 round-3 children unwalked; the deepening queue passed a canonical instead of the file's own orientation for 19 seeds (34 walks) and was stopped in its first round; two rounds of closure materialization crashed at load; seven output files of the alternative-seed walks were overwritten by same-named outputs (none of them was a lowering). The cache-history defect of \cref{sec:cachebug} affected 43 of 1{,}098 replays (relative to the campaign's own pure-mode replays). We report these because they bound what the numbers can support, and because two of them (the cache-history defect and the self-referential selection) were found only when we tried to reproduce or re-derive results.

\paragraph{Comparators and claims.}
The tables are moving and incomplete (\cref{tab:pitfalls}); a result is a triple of scheme, date and revisions, and ``below all listed'' does not mean ``new''. A comparison with the four tables alone overstated novelty: the limited check of the companion paper superseded 87 of the 795 results (32 of the 80 rebuilt ones with F2-only schemes that our own closure lacked as leaves) and matched 112 more by compositions of published schemes and by published rank claims. It was not a complete prior-art audit. Some of the schemes that seed our results are composed from constructions that the project's own notes record as established.

\paragraph{Scope.}
The pipeline finds upper bounds over $\F$ only. It gives no lower bounds and no information about lifting to characteristic zero, and nothing here bears on the exponent of matrix multiplication (the ranks are of small formats and the composition is standard).

\paragraph{Future work.}
The obvious steps are to walk the children that the defect left behind, to run matched-cost comparisons of uniform against directed sampling and of long walks against restarts, to promote the campaign drivers (currently scripts with scratch paths) into the package with parameters, to reduce the residual leaf-provenance gaps, to extend replay to a larger fraction of walks on a cluster, and to re-score against the moving tables at submission time.

\section{Conclusion}
A pipeline built from a directed flip-graph walker, a cheap cross-shape derivation, a rank-only composition closure, careful use of public schemes, and a verification protocol in which every result is checked twice and, wherever possible, re-derived from pinned upstream revisions, produced a large table of upper bounds over $\F$ (reported in the companion paper) at a cost of about $3\times10^{11}$ proposals.
The measured yields say where the effort paid: composition and projection supplied most of the records before the prior-art check, which then superseded or matched 128 of the 203 closure products against 35 of the 288 records with a walk or a projection; walks from already good public schemes supplied very few, and longer walks supplied almost none that the shorter ones had missed.
The engineering lessons are at least as general as the numbers: read the ring from the data and not from the directory; exclude commutative algorithms; treat every table as a dated, moving, field-agnostic bar; make every derivation replayable from a pinned revision; and verify from a clean extraction with empty caches, because a warm cache makes a verifier vacuous.

\section*{Availability}
The search and verification tools are in the package \texttt{tungsten-metaflip} (SPDX: MIT OR Apache-2.0 WITH LLVM-exception; bundled third-party data keeps its own licences, see \texttt{THIRD\_PARTY.md}) of the Tungsten repository \cite{tungsten2026lang}, branch \texttt{main} (the tools of the release are commit \texttt{00ffa04f9}; the sources of these papers, their scripts and the verification logs are in later commits) \todo{push \texttt{main} before submission}.
The tensors of the lineage bundles, the record index, the logs of the release verification and a snapshot of the tools are in a release of \texttt{tungsten-lang/metaflip-archives} (release \texttt{evidence-2026-09-29}), currently private \todo{deposit publicly, e.g.\ on Zenodo, and cite the DOI}.
Each bundle's README gives the exact commands that reproduce its verification; the log of the later Ruby check over all tensors and the closure rebuild log are in the repository (\texttt{papers/verification}).
The FastMatrixMultiplication repository is MIT-licensed, fmm-schemes is CC BY 4.0, the matmulcatalog carries no licence file, and the Lille catalogue is CC BY 3.0 according to its Wikidata record; schemes derived from them keep their attribution, recorded per row.
